\section{Choosing the parameters}\label{sec:parameters}

The choice of parameters has two constraints. The number of scalar
conditions in a query must be small compared with the ambient dimension
\(N=M_0n\). The number of channels must already suffice for matrix
completion before \(M_0\) is chosen. We give the choices in order.

\subsection{Pins and exclusions}

Begin with
\[
 D=8000,\qquad M=2^{1000},\qquad
 \zeta=\frac1{1000\cdot112}.
\]
The spreading estimate uses the allowance 112. A list in this
construction has at most 56 ambient key vectors by \eqref{eq:key-rank-budget}.

Closing the pins under projection leaves a gain of \(\zeta/6\) in the
exponent in Lemma~\ref{lem:closed-peeling}, before the raw image error.
Choose an initial
pin budget \(e\), a probe count \(k\), and then a second pin budget \(s\)
so that
\begin{align*}
 D+\zeta&\le\frac{\zeta}{8}e,& D+2&\le .49k,\\
 d&=4(2e+2)k,& u&=10(e+d+1),\\
 D+e+2+u+\zeta&\le\frac{\zeta}{8}s.&
\end{align*}
Here \(d\) bounds the preliminary cover size, and \(u\) allows the
additional pins for the marks and cover actions. The final pin budget is
\(K=e+s\). Rounding each lower bound upward gives integers.
All these choices depend only on the first displayed constants.
Now choose the odd cycle length
\[
 g=4e+20001.
\]
The cycle has degree two regardless of its length. Thus increasing \(g\)
does not change the slot allowance or the early pin budgets.

Put
\[
 R=4(K+28),\qquad q_*=R^2+28,\qquad
 d_S=2gKq_*,\qquad H_*=2^{d_S+34}.
\]
There are two endpoint primal pin spaces at each of the \(g\) components.
Thus \(d_S\) bounds the protected selector dimension in
\eqref{eq:label-budget}. The allowance \(H_*\) covers any selector span
formed while choosing the short lists. Polynomial continuation preserves
the spans throughout these choices. We retain the conservative exclusion
budget
\[
 B_{\rm exc}=10(10H_*+10000g+1).
\]
This includes the prediction exclusions and all scalar selection losses.
Now choose \(b\) so that
\[
 2^{b-2}>100g(B_{\rm exc}+1).
\]
In particular \eqref{eq:selector-size} holds. The selector degree is one
and its dimension is \(p_*=b+1\). The protected base dimension is at most \(2gKp_*\). This bound is
used when \(n\) is chosen in \eqref{eq:base-rejection}.

\subsection{Mixers and channels}

The count of profiles of total rank at most \(K\) can use
\[
 A=10(K+1)p_*(g^2+5),\qquad s_0=4A+1.
\]
Choose
\[
 J=100(s_0+K^2+1)+1.
\]
These bounds give the unary mixing union bound of
Lemma~\ref{lem:generic-mixers}. The explicit binary construction has \(2Jg\)
indexed matrices. Its shift degrees are bounded before \(n\), so their
required ranks hold for all sufficiently large \(n\).

Choose \(r_0\) next, always setting \(h=1000r_0\). We check that increasing
\(r_0\) meets the two rank demands that also involve \(h\).

For realization, the conservative joint gradient rank is
\[
 r_{\rm grad}=2(15r_0+14Jg).
\]
For \(t\le28\), put \(f=2(K+t)\). The residual quadratic compression
in \Cref{sec:realization} has a bound \(d_{\rm cmp}(f,3f+K)\) that depends
only on \(g,p_*,f,K\). One may take
\[
 \begin{split}
 d_{\rm bits}(T,R')&=2g p_*(T+R'),\\
 d_{\rm row}(T,R')&=p_*+p_*(g^2+3)d_{\rm bits}(T,R'),\\
 d_{\rm cmp}(f,r)&=2d_{\rm row}(2f,2r).
 \end{split}
\]
The extension requires
\begin{equation}\label{eq:ledger-realization}
 r_{\rm grad}+f+K+d_{\rm cmp}(f,3f+K)+5K+6f+2\le h.
\end{equation}
The left side is \(30r_0+C\), where \(C\) is independent of \(r_0,h,n\).
It suffices to take \(970r_0\ge C\), using the largest value at \(t=28\).
The coefficient 30 is independent of the pin dimension because the
tester contribution was removed before compressing the residual.

For the prediction identity the rank bound on the terms outside the
ordinary and shared testers is
\[
 r_{\rm scr}=2(2h+2Je).
\]
The exact support argument requires
\begin{equation}\label{eq:ledger-tester}
 4r_{\rm scr}+2er_0<gr_0.
\end{equation}
After substituting \(h=1000r_0\), this is
\[
 (16000+2e)r_0+16Je<gr_0.
\]
Our choice of \(g\) makes this hold whenever \(r_0\ge Je+1\).
The scalar extraction also uses marks in the later pin space, of bound
\(K=e+s\). It only needs
\[
 4\cdot2(2h+2JK)<gr_0.
\]
Since \(g>20000\), this follows after also requiring \(r_0\ge JK+1\).
Thus both requirements can be met by increasing \(r_0\) after
\(J,K\) are fixed.

The phase and safety estimates also impose lower bounds on \(h\).
Their cover sizes, probe counts, and error tolerances were fixed using
\(D,e,s,u,k\). For example, the moment calculation compares
\(4^a p_0^{qa/2}\) and \(2^{-hqa/4}\) with a fixed exponential density
cost, where \(q=2^{-100r-10}\) and \(a\) is a small fixed multiple of \(N\).
Choose \(p_0\) small using the early constants, then choose \(h\) large.
For \(k\) probe directions, the bound in the proof of
Proposition~\ref{thm:phase-alternative} has exponent
\((D+1+.012-.49k)N\). The choice \(.49k\ge D+2\) makes this at
most \(-.988N\). The remaining failures have factors
such as \(2^{L-h}\), with \(L\) already fixed. They too become as small
as required. No common key partition or later accuracy enters these
choices.

\subsection{The ambient dimension}

The nominal dimension is
\[
 d_n=p_*\bigl(1+(g^2+3)n\bigr).
\]
Choose an even \(M_0\) sufficiently large and put \(N=M_0n\).
Then \(d_n/N\) and every other fixed multiple of \(n/N\) are as small
as required. In particular, there is room for the injective quadratic
frames and their unused channels. The exact-image estimates can use an
error smaller than \(\zeta/24\).

The records in the collision proof have \(2^{O(n)}\) possible values.
Choose \(M_0\) so that their logarithm is at most \(.005N\). The same
choice makes \(16ghd_n\le .01N\), as required for the scalar entries
tested in Proposition~\ref{prop:collision}. All coefficients in these expressions
were fixed before \(M_0\).

The resulting order is
\[
 (D,M,\zeta)\ \longrightarrow\ (e,k,d,u,s,K)
 \ \longrightarrow\ (g,B_{\rm exc})
 \ \longrightarrow\ (b,s_0,J)
 \ \longrightarrow\ (r_0,h)
 \ \longrightarrow\ M_0\ \longrightarrow\ n.
\]
The notation \(B_{\rm exc}\) in this line includes the intermediate
quantities \(R,q_*,d_S,H_*\), all determined by \(g,K\).

\subsection{Later analytic choices}

The Fourier comparison in \Cref{sec:histograms} begins with a fixed
leaf density cap and a desired error. Once the repetition rate is chosen,
fix the cutoff for nominal rank deficiency. The cutoff for cross rank
can then absorb the cost of conditioning on the nominal kernel.
Finally choose \(\delta\) for the low rank contribution in
\eqref{eq:assembled-moment}. This determines how large \(n\) must be.

For the linear partitions the bound \(V\) in
\eqref{eq:linear-density-energy} is independent of their number of rows.
Lemma~\ref{lem:linear-test} uses Parseval to bound the number of
frequencies needed at a prescribed accuracy. The height \(2B_0\) used
for truncation is also independent of that accuracy. In the scalar
argument the overlap threshold is fixed below the gap in
\eqref{eq:status-support-gap} before choosing the common partition.
In the case with predictions, every attained predictor cell already
has mass at least \(2^{-2R_*}/2\).

\subsection{Final probability margins}

A unit query has at most 28 atom positions. Each atom uses two cycle
components and contributes one ambient key vector on each. Thus there
are at most 56 vectors. The spreading loss satisfies
\[
 2\zeta\cdot112\le .004.
\]
The discarded records and the tests of scalar entries in
Proposition~\ref{prop:collision} give the final exponent \(56+.03\). The retained mass has a fixed positive lower
bound. An overlap of order \(N^{-4}\) costs only \(O(\log N)\) in the
exponent. Hence
\[
 2^{-(56+.03)N-o(N)}>2^{-200N}
\]
for large \(n\).

Finally take \(m=2^{2000N}\), \(\varepsilon=2^{-200N}\), and density
cap \(\kappa=2^{8000N}\) in the graph transfer. Each sample records a list of
\(2^{400N}\) matching edges. Since \(\log|\Omega|=O(N^2)\), the logarithm of the
number of lists is \(O(N^2 2^{400N})=o(m)\).
The two sampling errors are exponentially small in \(m\).
